%% ==================================================
%% Chapter: Sequential generation of matrix product states
%% Source: arXiv:quant-ph/0608197, Section 5 (lines 1504--1616 of
%% Papers/quant-ph_0608197/MPSarchive.tex); arXiv:quant-ph/0501096.
%% ==================================================

\chapter{Sequential Generation of Matrix Product States}\label{ch:sequential_generation}

A chain of $N$ sites with local dimension $d$ starts in the product state
$\ket{0}^{\otimes N}$. In a sequential scheme an ancilla of dimension $D$
interacts with the sites one after another, or the sites interact with their
neighbours one after another. The states produced in this way are exactly the
matrix product states with open boundary conditions and bounded bond
dimension. The presentation follows \cite[Section~5]{PerezGarcia2007Matrix},
which reviews and extends \cite{Schoen2005Sequential}.

\section{Sequential schemes}

\begin{definition}[Open-boundary matrix product state]
    \label{def:obc_mps_bond_bound}
    A vector $\ket{\psi}\in(\C^d)^{\otimes N}$ has an \emph{open-boundary matrix
        product representation with bond dimension at most $D$} if there are
    integers $D_1,\dots,D_{N+1}\le D$ with $D_1=D_{N+1}=1$ and matrices
    $A^{[k]}_i\in\C^{D_{k+1}\times D_k}$, for $k=1,\dots,N$ and $i=0,\dots,d-1$,
    such that
    \begin{align}
        \ket{\psi}
         & =\sum_{i_1,\dots,i_N=0}^{d-1}
        A^{[N]}_{i_N}\cdots A^{[2]}_{i_2}A^{[1]}_{i_1}\,
        \ket{i_N\cdots i_1}.
        \label{eq:obc_mps_bond_bound}
    \end{align}
    Here the product of the $1\times 1$ matrices is read as a scalar.
\end{definition}

\begin{definition}[Sequential scheme with an ancilla]
    \label{def:sequential_scheme_ancilla}
    \uses{def:obc_mps_bond_bound}
    Let $\ket{\varphi_I},\ket{\varphi_F}\in\C^D$. For $k=1,\dots,N$ let
    $G_k=\sum_{i,\alpha,\beta}A^{[k]}_{i,\alpha,\beta}\,
        \ketbra{\alpha,i}{\beta,0}$ act on the ancilla $\C^D$ and the $k$-th site.
    Applying $G_1,\dots,G_N$ in turn to $\ket{\varphi_I}\otimes\ket{0}^{\otimes N}$
    and projecting the ancilla onto $\ket{\varphi_F}$ leaves the chain in
    \begin{align}
        \ket{\psi}
         & =\sum_{i_1,\dots,i_N=0}^{d-1}
        \bra{\varphi_F}A^{[N]}_{i_N}\cdots A^{[2]}_{i_2}A^{[1]}_{i_1}
        \ket{\varphi_I}\,\ket{i_N\cdots i_1},
        \label{eq:sequential_generated_state}
    \end{align}
    where $(A^{[k]}_i)_{\alpha\beta}=A^{[k]}_{i,\alpha,\beta}$; here
    $\alpha,\beta$ index the ancilla, so each $A^{[k]}_i$ is a $D\times D$
    matrix.
    The scheme is
    \begin{enumerate}
        \item \emph{probabilistic} if the operations $G_k$ are arbitrary;
        \item \emph{deterministic} if every $G_k$ is the restriction of a unitary,
              so that $\sum_i (A^{[k]}_i)^\dagger A^{[k]}_i=\Id_D$, and the ancilla
              decouples after the last step, that is, the final joint state is
              $\ket{\varphi_F}\otimes\ket{\psi}$;
        \item for $d=2$, a \emph{deterministic transition scheme} if the ancilla is
              $\C^D\otimes\C^2$, every step consists of an arbitrary unitary on the
              ancilla followed by the fixed interaction
              $\ket{\varphi}\ket{1}\ket{0}\mapsto\ket{\varphi}\ket{0}\ket{1}$,
              $\ket{\varphi}\ket{0}\ket{0}\mapsto\ket{\varphi}\ket{0}\ket{0}$, and the
              ancilla decouples after the last step.
    \end{enumerate}
    This is \cite[Section~5.1]{PerezGarcia2007Matrix}.
\end{definition}

\section{Sequential generation with an ancilla}

\begin{lemma}[Isometric open-boundary representation]
    \label{lem:obc_left_isometric_representation}
    \uses{def:obc_mps_bond_bound}
    Every normalized vector with an open-boundary matrix product representation
    with bond dimension at most $D$ has such a representation with the same
    bound in which every site map is an isometry, that is, for $k=1,\dots,N$,
    \begin{align}
        \sum_{i=0}^{d-1}(A^{[k]}_i)^\dagger A^{[k]}_i
         & =\Id_{D_k}.
        \label{eq:obc_left_isometric}
    \end{align}
\end{lemma}
\begin{proof}
    Starting from the last site $k=N$, read the site map as a matrix from
    $\C^{D_k}$ to $\C^{D_{k+1}}\otimes\C^d$ and write it as $Q_kR_k$ with
    $R_k\colon\C^{D_k}\to\C^{r_k}$, $r_k\le D_k$, and $Q_k$ an isometry, for
    instance by a singular value decomposition. Keep $Q_k$ as the new site map,
    replace $D_k$ by $r_k$, and multiply $R_k$ into the site map of site $k-1$.
    The bond dimensions do not grow. After the first
    site the remaining factor is a $1\times 1$ matrix, and it has modulus one
    because the vector is normalized; absorb it into the first site map. This is the construction of \cite{Schoen2005Sequential} by
    successive decompositions.
\end{proof}

\begin{lemma}[Unitary completion of an isometry]
    \label{lem:isometry_unitary_completion}
    Let $W$ be a finite-dimensional Hilbert space, $S\subseteq W$ a subspace,
    and $V\colon S\to W$ an isometry. Then there is a unitary $U$ on $W$ with
    $U|_S=V$.
\end{lemma}

\begin{theorem}[Sequential generation with an ancilla]
    \label{thm:sequential_generation_ancilla}
    \uses{def:sequential_scheme_ancilla, def:obc_mps_bond_bound}
    The following sets of vectors in $(\C^d)^{\otimes N}$ coincide:
    \begin{enumerate}
        \item the states produced by probabilistic sequential schemes with a
              $D$-dimensional ancilla, up to normalization;
        \item the states produced by deterministic sequential schemes with a
              $D$-dimensional ancilla;
        \item for $d=2$, the states produced by deterministic transition schemes;
        \item the states with an open-boundary matrix product representation
              with bond dimension at most $D$, up to normalization.
    \end{enumerate}
    This is \cite[Theorem~13, Section~5.1]{PerezGarcia2007Matrix}.
\end{theorem}
\begin{proof}
    \uses{lem:obc_left_isometric_representation,
        lem:isometry_unitary_completion}
    Every generated state has the form (\ref{eq:sequential_generated_state}),
    which is an open-boundary representation with bond dimension at most $D$
    after absorbing $\ket{\varphi_I}$ and $\bra{\varphi_F}$ into the first and
    last site. Conversely, bring a normalized open-boundary state into the
    isometric form of Lemma~\ref{lem:obc_left_isometric_representation}. Each
    site map is an isometry from the bond space $\C^{D_k}$, embedded in the
    ancilla as $\C^{D_k}\otimes\ket{0}$, into $\C^{D_{k+1}}\otimes\C^d$;
    Lemma~\ref{lem:isometry_unitary_completion} extends it to a unitary on
    $\C^D\otimes\C^d$. Since $D_{N+1}=1$, the ancilla is left in a fixed
    state after the last step.
\end{proof}

\section{Sequential generation without an ancilla}

\begin{definition}[Sequential scheme without an ancilla]
    \label{def:sequential_scheme_no_ancilla}
    Starting from $\ket{0}^{\otimes N}\in(\C^d)^{\otimes N}$, apply an
    operation $U^{[k]}$ to sites $k$ and $k+1$ for $k=1,\dots,N-1$ in turn. The
    scheme is \emph{deterministic} if every $U^{[k]}$ is unitary and
    \emph{probabilistic} if the operations are arbitrary.
\end{definition}

\begin{theorem}[Sequential generation without an ancilla]
    \label{thm:sequential_generation_no_ancilla}
    \uses{def:sequential_scheme_no_ancilla, def:obc_mps_bond_bound}
    The states produced by sequential schemes without an ancilla, either
    deterministic or probabilistic (the latter up to normalization), are
    exactly the states with an open-boundary matrix product representation
    with bond dimension at most $d$. This is
    \cite[Theorem~14, Section~5.2]{PerezGarcia2007Matrix}.
\end{theorem}
\begin{proof}
    \uses{thm:sequential_generation_ancilla}
    For $k<N-1$ the matrices $A^{[k]}_{i,\alpha,\beta}=\bra{i,\beta}U^{[k]}\ket{\alpha,0}$
    give the site maps, with $\alpha=0$ at $k=1$; a singular value decomposition
    of the coefficients $\bra{i,j}U^{[N-1]}\ket{\alpha,0}$ in the grouping
    $i\mid\alpha,j$ gives the last two sites. Conversely, run the deterministic
    scheme of Theorem~\ref{thm:sequential_generation_ancilla} using site $k+1$
    as the ancilla of the $k$-th step, followed, for $k\le N-2$, by a swap of
    sites $k+1$ and $k+2$. The steps $k=1,\dots,N-1$ suffice: the last step
    of the deterministic scheme is a swap of the ancilla with site $N$, and at
    step $N-1$ the ancilla already is site $N$, so no operation is needed
    there \cite[proof of Theorem~14]{PerezGarcia2007Matrix}.
\end{proof}
